\section{数论}
\subsection{欧拉筛}
\subsubsection{线性筛质数}
\begin{lstlisting}[caption={}]
bitset<N> isprime;
vector<i64> prime;
void euler() {
    isprime.set();
    isprime[1] = 0;
    for (i64 i = 2; i <= N; i++) {
        if (isprime[i]) prime.push_back(i);
        for (int j = 0; j < prime.size() && i * prime[j] < N; j++) {
            isprime[i * prime[j]] = false;
            if (i % prime[j] == 0) break;
        }
    }
}
\end{lstlisting}

\subsubsection{线性筛因数个数}
\begin{lstlisting}[caption={}]
vector<int> d(N);      // 因数个数
vector<int> minp_cnt(N, 1);  // 记录最小质因子的指数+1
void euler_d(){
    d[1] = 1;
    for(ll i = 2; i < N; i++){
        if(isprime[i]){
            d[i] = minp_cnt[i] = 2;
            prime.push_back(i);
        }
        for(int j = 0; j < prime.size() && i * prime[j] < N; j++){
            i64 p = prime[j];
            if(i % p == 0){
                minp_cnt[i * p] = minp_cnt[i] + 1;
                d[i * p] = d[i] / minp_cnt[i] * minp_cnt[i * p];
                break;
            }
            else{
                d[i * p] = d[i] * 2;
                minp_cnt[i * p] = 2;
            }
        }
    }
}
\end{lstlisting}

\subsubsection{线性筛因数和}
\begin{lstlisting}[caption={}]
vector<i64> sigma(N);  // 因数和
vector<i64> minp_pow(N, 1);  // 记录最小质因子的幂次和
void euler_sigma(){
    sigma[1] = 1;
    for (i64 i = 2; i < N; i++) {
        if(isprime[i]){
            sigma[i] = minp_pow[i] = i + 1;
            prime.push_back(i);
        }
        for(int j = 0; j < prime.size() && i * prime[j] < N; j++){
            i64 p = prime[j];
            if(i % p == 0){
                minp_pow[i * p] = minp_pow[i] * p + 1;
                sigma[i * p] = sigma[i] / minp_pow[i] * minp_pow[i * p];
                break;
            }
            else{
                sigma[i * p] = sigma[i] * sigma[p];
                minp_pow[i * p] = 1 + p;
            }
        }
    }
}
\end{lstlisting}

\subsection{逆元}
如果一个线性同余方程 $a^x \equiv 1 \pmod b$，则 $x$称为$a \bmod b$ 的逆元，记作$ a^{-1}$。
\subsubsection{费马小定理}
设 $p$ 是素数。对于任意整数 $a$ 且$p \perp a$,都成立$a^{p-1}\equiv 1\pmod p$.
\begin{lstlisting}[caption={}]
i64 inv(i64 x){
    return qpow(x,mod - 2);
}
\end{lstlisting}

\subsubsection{扩展欧几里得算法}
其中 $x$ 为 $a$ 的逆元
\begin{lstlisting}[caption={}]
int exgcd(int a,int b,int &x,int &y){
	if(b == 0){
		x = 1;
		y = 0;
		return a;
	}
	int d = exgcd(b,a % b,y,x);
	y -= (a / b) * x;
	return d;
}
\end{lstlisting}

\subsubsection{线性求乘法逆元}
$$i^{-1} \equiv - \left\lfloor \dfrac{p}{i} \right\rfloor (p\bmod i)^{-1} \pmod p$$
\begin{lstlisting}[caption={}]
inv[1] = 1;
for(int i = 2;i <= n;i++)
	inv[i] = (long long)(p - p / i) * inv[p % i] % p;
\end{lstlisting}

\subsection{扩展大步小步算法/exBSGS}
以 $\mathcal O(\sqrt {P})$ 的复杂度求解 $a^x \equiv b \pmod P$ 的最小自然数 $x$。
其中标准 $BSGS$ 算法不能计算 $a$ 与 $MOD$ 互质的情况，而 $exbsgs$ 则可以。
\begin{lstlisting}[caption={}]
i64 BSGS(i64 a,i64 b,i64 p,i64 k = 1){
    i64 A = sqrtl(p) + 1;
    map<i64,i64> mp; 
    i64 cur = 1;
    for(int i = 1;i <= A;i++){
        cur = (cur * a) % p;
        mp[(cur * b) % p] = i;
    }
    i64 temp = cur * k % p;
    for(int i = 1;i <= A;i++){
        auto it = mp.find(temp);
        if(it != mp.end()) return i * A - it->second;
        temp = (temp * cur) % p;
    }
    return -1e18;
}
i64 exBSGS(i64 a, i64 b, i64 m, i64 k = 1){  //O(/sqrt(m)logm)
    a %= m,b %= m;
    i64 A = a, B = b, M = m;
    if (b == 1) return 0;
    i64 cur = 1 % m;
    for (int i = 0;; i++){
        if (cur == B) return i;
        cur = cur * A % M;
        i64 d = __gcd(a, m);
        if (b % d) return -1e18;
        if (d == 1) return BSGS(a, b, m, k * a % m) + i + 1;
        k = k * a / d % m, b /= d, m /= d;
    }
}
int main(){
    ios::sync_with_stdio(0),cin.tie(0),cout.tie(0);
    i64 a,b,p;
    cin >> a >> p >> b;
    i64 temp = exBSGS(a,b,p);
    if(temp >= 0) cout << temp << "\n";
    else cout << "No Solution\n";
    return 0;
}
\end{lstlisting}

\subsection{欧拉函数}
若 $n = p^k$，其中 $p$ 是质数，那么 $\varphi(n) = p^k - p^{k - 1}$

由唯一分解定理，设 
$n = \prod_{i=1}^{s}p_i^{k_i}$，其中$p_i$ 是质数，有 
$\varphi(n) = n \times \prod_{i = 1}^s{\dfrac{p_i - 1}{p_i}}$。
\subsubsection{单个数的欧拉函数}
$1$ 到 $N$ 中与 $N$ 互质数的个数称为欧拉函数，记作 $\varphi (N)$ 。
求解欧拉函数的过程即为分解质因数的过程，复杂度 $\mathcal{O}(\sqrt{n})$ 。
\begin{lstlisting}[caption={}]
int euler_phi(int n) {
    int ans = n;
    for (int i = 2; i * i <= n; i++)
        if (n % i == 0) {
        ans = ans / i * (i - 1);
        while (n % i == 0) n /= i;
        }
    if (n > 1) ans = ans / n * (n - 1);
    return ans;
}
\end{lstlisting}

\subsubsection{线性筛欧拉函数}
\begin{lstlisting}[caption={}]
vector<int> phi(N);
void euler_phi(){
    phi[1] = 1;
    for(i64 i = 2; i < N; i++){
        if(isprime[i]){
            phi[i] = i - 1;
            prime.push_back(i);
        }
        for(int j = 0; j < prime.size() && i * prime[j] < N;j++) {
            i64 p = prime[j];
            if(i % p == 0){
                phi[i * p] = phi[i] * p;
                break;
            }
            else phi[i * p] = phi[i] * (p - 1);
        }
    }
}
\end{lstlisting}

\subsubsection{欧拉函数常见性质}
\begin{itemize}
    \item $1-n$ 中与 $n$ 互质的数之和为 $n \cdot \frac{\varphi(n)}{2}$;
    \item 若 $a,b$ 互质，则 $\varphi (a \cdot b) = \varphi (a)  \cdot \varphi(b)$ 。实际上，所有满足这一条件的函数统称为积性函数;
    \item 若 $f$ 是积性函数，且有 $\displaystyle n = \prod ^m _{i =1} p_i ^ {c_i}$ ，那么 $\displaystyle f(n) = \prod ^m _{i =1} f( p_i ^ {c_i} )$;
    
    \item 若 $p$ 为质数，且满足 $p \mid  n$:
    \begin{itemize}
        \item $n \bmod p^2 = 0$ ，那么 $\varphi (n) = \varphi (\frac{n}{p})  \cdot p$ ;
        \item $n \bmod p^2 \neq 0$，那么 $\varphi (n) = \varphi (\frac{n}{p})  \cdot (p-1)$ ;
    \end{itemize}
    \item 欧拉反演: $\displaystyle\sum _{d \mid n} \varphi (d)= n$ ;

          如 $n=10$ ，则 $d=10/5/2/1$ ，那么 $10 = \varphi(10) + \varphi(5) + \varphi(2) + \varphi(1)$ ;
    \item $\displaystyle\sum_{i = 1}^{n} \gcd(i, n) = \sum_{d|n} \left\lfloor \frac{n}{d} \right\rfloor \varphi(d)$;
    \item $\varphi(n) = \displaystyle\sum_{d | n} \mu(d) \frac{n}{d}$;
    \item $\displaystyle\sum_{i = 1}^n\varphi(n) = \displaystyle\sum_{d = 1}^n \mu(d) \frac{\left\lfloor \frac{n}{d} \right\rfloor(\left\lfloor \frac{n}{d} \right\rfloor + 1)}{2}$;
\end{itemize}


\subsection{莫比乌斯函数}
$\mu$ 为莫比乌斯函数，定义为:
        $$\mu(n)= \begin{cases}
        1&n=1\\
        0&n\text{ 含有平方因子}\\
        (-1)^k&k\text{ 为 }n\text{ 的本质不同质因子个数}\\
        \end{cases}$$
\subsubsection{单个数的莫比乌斯函函数}
 $\mathcal{O}(\sqrt{n})$ 。
\begin{lstlisting}[caption={}]
int mu(int n) {
    int res = 1;
    for(int i = 2; i * i <= n;i++){
        if(n % i == 0){
            n /= i;
            if (n % i == 0) return 0;
            res = -res;
        }
    }
    if (n > 1) res = -res;
    return res;
}
\end{lstlisting}

\subsubsection{线性筛莫比乌斯函数}
\begin{lstlisting}[caption={}]
vector<int> mu(N);
void euler_mu(){
    mu[1] = 1;
    for(i64 i = 2; i < N; i++){
        if(isprime[i]){
            mu[i] = -1;
            prime.push_back(i);
        }
        for(int j = 0;j < prime.size() && i * prime[j] < N; j++){
            i64 p = prime[j];
            if(i % p == 0){
                mu[i * p] = 0;
                break;
            }
            else mu[i * p] = -mu[i];
        }
    }
}
\end{lstlisting}

\subsubsection{莫比乌斯常见性质}
\begin{itemize}
    
    \item 常见积性函数之间的关系:
        \begin{itemize}
            \item 常数函数: $1(n) = 1$
            \item 单位函数: $\epsilon(n) = [n = 1] = \displaystyle\sum_{d | n}\mu(d) \Leftrightarrow \epsilon = \mu \ast 1 $;$f * \varepsilon = \varepsilon * f = f$
            \item 恒等函数: $id(n) = n = \displaystyle\sum _{d \mid n} \varphi (d) \Leftrightarrow \displaystyle id = \varphi \ast 1$
            \item 幂  函数: $id_k(n) = n^k$
            \item 约数个数: $\tau(n) = \displaystyle\sum_{d | n}1  \Leftrightarrow \tau = 1 \ast 1$
            \item 约数和  : $\sigma(n) = \displaystyle\sum_{d | n}d  \Leftrightarrow \sigma = id \ast 1$
            \item 约数函数: $\sigma_k(n) = \displaystyle\sum_{d | n}d^k  \Leftrightarrow \sigma_k = id_k \ast 1$
            \item 欧拉函数: $\varphi(n) = \displaystyle\sum_{i = 1}^x[gcd(x,i) = 1] = \displaystyle\sum_{d | n} d \cdot \mu(\frac{n}{d}) \Leftrightarrow \varphi = id \ast \mu$
            \item 莫比乌斯函数: $\mu = \mu \ast \epsilon$
        \end{itemize}

    \item $[gcd(a,b) == 1] = \sum\limits_{d | gcd(a,b)} \mu(d)$
\end{itemize}
       
\subsubsection{莫比乌斯反演}
如果有 $f(n)=\sum\limits_{d\mid n}g(d)$，那么有 $g(n)=\sum\limits_{d\mid n}\mu(d)f(\frac{n}{d})$。

如果 $g = f * 1$，那么 $f = g * \mu$

一般流程:求解 $\displaystyle \sum^{n}_{i = 1}\sum^{m}_{j = 1}f(gcd(i,j))$
\begin{itemize}
    \item 考虑构造出函数$g$,满足下列形式：
            $$f(n) = g \ast 1 = \displaystyle\sum_{d|n} g(d)$$
    \item 则求和式子变为:
        $$\displaystyle \sum^{n}_{i = 1}\sum^{m}_{j = 1}\sum_{d|gcd(i,j)} g(d)$$
    \item 考虑算术基本定理，发现若满足$d|gcd(i,j)$，则 $d|i $且 $d|j$ 成立。考虑 $g(d)$在何时做出贡献，调整枚举顺序：
    $$\displaystyle \sum_{d = 1}g(d) \sum^{n}_{i = 1} [d | i] \sum^{m}_{j = 1} [d | j]$$

    \item $\sum^{n}_{i = 1} [d | i]$等价于$1\sim n$中$d$的倍数的个数,则上式等价于:
     $$\displaystyle \sum_{d = 1}g(d) \left \lfloor \frac{n}{d} \right \rfloor \left \lfloor \frac{m}{d} \right \rfloor$$

\end{itemize}
    
\subsection{数论分块}
$\displaystyle \left\lfloor \frac{n}{l} \right\rfloor = \left\lfloor \frac{n}{l + 1} \right\rfloor
 = ... = \left\lfloor \frac{n}{r} \right\rfloor \iff \left\lfloor \frac{n}{l} \right\rfloor \le \frac{n}{r} 
 < \left\lfloor \frac{n}{l} \right\rfloor + 1$ ，根据不等式左侧，
 得到 $\displaystyle r \le \left\lfloor \frac{n}{\lfloor \frac{n}{l} \rfloor} \right\rfloor$ 。
\begin{lstlisting}[caption={}]
void solve(i64 n){
    i64 ans = 0;
    for(i64 i = 1, j; i <= n; i = j + 1){
        j = n / (n / i);
        ans += (j - i + 1) * (n / i);
    }
    cout << ans << "\n";
}
\end{lstlisting}

\subsection{矩阵快速幂}
\begin{lstlisting}[caption={}]
static constexpr int mod = 1e9 + 7,N = 2;
using Martix = array<array<int,N>,N>;
Martix mul(Martix& a,Martix& b){
    Martix res{};
    for(int k = 0;k < N;k++){
        for(int i = 0;i < N;i++){
            if(a[i][k] == 0) continue;
            for(int j = 0;j < N;j++){
                res[i][j] = (res[i][j] + ((long long)a[i][k] * b[k][j]) % mod) % mod;
            }
        }
    }
    return res;
}
Martix qpow(Martix& a,int b){
    Martix temp{};
    for(int i = 0;i < N;i++) temp[i][i] = 1;
    while(b){
        if(b & 1) temp = mul(temp,a);
        a = mul(a,a);
        b >>= 1;
    }
    return temp;
}
\end{lstlisting}

\subsection{异或空间线性基(贪心)}
\begin{lstlisting}[caption={}]
struct xorBase{
    vector<i64> p;
    xorBase():p(61){}
    void insert(i64 x){
        for(int i = 60;i >= 0;i--)
            if(x >> i){
                if(p[i]) x ^= p[i];
                else{p[i] = x;return;}
            }
    }

    i64 getmx() {
        i64 ans = 0;
        for(int i = 60;i >= 0;i--)
            if((ans ^ p[i]) > ans) ans ^= p[i];
        return ans;
    }
};
\end{lstlisting}
\subsubsection{常见操作}
\begin{itemize}
    \item 求最小值：如果是求用线性基 $p$内的元素能异或出的最小值，那么就是最小的$p_i$。
                如果是求整个序列能异或出的最小值的话，需检查有没有元素不能插入线性基，如果有，那么最小值就是$0$，否则依然是最小的$p_i$​。
    \item 求$k$小值:从一个序列中取任意元素进行异或，求能异或出的所有数字中第$k$小的那个。
                首先处理这个序列的线性基 $p$，对于每一个 $p_i$,枚举$j=i\sim 1$,如果$p_i$ 二进制的第 $j$ 位为 $1$，那么 $p_i$ 异或上 
                $p_{j-1}$。将 $k$ 先转成二进制，假如第 $i$  位为 $1$,则 $ans$ 异或线性基中第 $i$个元素(注意不是直接异或 $p_{i-1})$。
    \item 询问存在性:判断数 x 能否被当前线性基中元素异或得到。把 $x$ 尝试插入进线性基，假如可以插入，说明不能异或得到，假如插不进去，则说明可以异或得到。
\end{itemize}




\subsection{组合数}
\subsubsection{通项公式}
\begin{lstlisting}[caption={}]
i64 fac[N],invfac[N];
i64 calc(i64 m,i64 n){
    return (fac[n] * ((invfac[m] * invfac[n - m]) % mod)) % mod;
}
void init(){
    fac[0] = fac[1] = 1;
    invfac[0] = invfac[1] = 1;
	for(int i = 2;i < N;i++){
		fac[i] = i;
		invfac[i] = (mod - mod / i) * invfac[mod % i] % mod;
	}
    for(int i = 2;i < N;i++){
		fac[i] = fac[i] * fac[i - 1] % mod;
		invfac[i] = invfac[i] * invfac[i - 1] % mod;
    }
}
\end{lstlisting}

\subsubsection{卢卡斯定理}
$$ C^n_m \equiv C^{n \bmod p}_{m \bmod p} \cdot C^{\left \lfloor n/p  \right \rfloor}
_{\left \lfloor m/p  \right \rfloor} \pmod p $$
\begin{lstlisting}[caption={}]
i64 C(i64 m, i64 n, i64 p){
    return m < n ? 0 : fact[m] * inv(fact[n], p) % p * inv(fact[m - n], p) % p;
}
i64 lucas(i64 m, i64 n, i64 p){
    return n == 0 ? 1 % p : lucas(m / p, n / p, p) * C(m % p, n % p, p) % p;
}
\end{lstlisting}

\subsection{中国剩余定理}
\begin{lstlisting}[caption={}]
template<typename T>
T exgcd(T a,T b,T &x,T &y){
	if(b == 0){
		x = 1;
		y = 0;
		return a;
	}
	T d = exgcd(b,a % b,y,x);
	y = y - (a / b) * x;
	return d;
}

template<typename T>
T inv(T a,T p){
    T x,y;
    exgcd(a,p,x,y);
    return (x % p + p) % p;
}

int main(){
    ios::sync_with_stdio(0),cin.tie(0),cout.tie(0);
    int T = 1;
    // cin >> T;
    while(T--){
        ll n;
        __int128 M = 1;
        cin >> n;
        vector<ll> a(n + 1),m(n + 1);
        for(int i = 1;i <= n;i++){
            cin >> m[i] >> a[i];
            M *= m[i];
        }
        __int128 sum = 0;
        for(int i = 1;i <= n;i++){
            __int128 tm = M / m[i];
            sum = (sum + ((__int128)a[i] * tm) % M * inv(tm,(__int128)m[i]) % M) % M;
        }
        cout << (ll)sum << "\n";
    }   
    return 0;
}
\end{lstlisting}

\subsection{快速傅里叶变换(FFT)}
\begin{lstlisting}[caption={}]
#include<bits/stdc++.h>
using namespace std;

typedef long long ll;
const int N = 4e6 + 10;
const double pi = acos(-1.0);

complex<double> f[N], g[N];
int R[N];
void change(complex<double> A[],ll n){
	for(int i = 0;i < n;i++)
		R[i] = R[i / 2] / 2 + (i & 1) * (n / 2);
	for(int i = 0;i < n;i++)
		if(i < R[i]) swap(A[i],A[R[i]]);
}

void FFT(complex<double> A[],ll n,int op){ //op为虚部符号，op为1时FFT，op为-1时IFFT
	change(A,n); //位逆序变换
	for(int m = 2;m <= n;m <<= 1){  //枚举位宽 
		complex<double> w1({cos(2 * pi / m),sin(2 * pi / m) * op});
		for(int i = 0;i < n;i += m){  //枚举块数 
			complex<double> wk({1,0});
			for(int j = 0;j < m / 2;j++){
				complex<double> x = A[i + j],y = A[i + j + m / 2] * wk;
				A[i + j] = x + y;
				A[i + j + m / 2] = x - y;
				wk = wk * w1;
			}
		}
	} 
}

int main(){
	ios::sync_with_stdio(0),cin.tie(0),cout.tie(0);
	ll n,m;
	cin >> n >> m;
	for(ll i = 0;i <= n;i++) cin >> f[i];
	for(ll i = 0;i <= m;i++) cin >> g[i];
	ll len = 1 << max((ll)ceil(log2(n + m)),1LL);  //FFT需要项数为2的整数次方倍,len为第一个大于a.size() + b.size()的二的正整数次方
	FFT(f,len,1),FFT(g,len,1);  //系数表达转点值表达
	for(int i = 0;i <= len;i++)
		f[i] = f[i] * g[i];
	FFT(f,len,-1);  //点值表达转系数表达
	for(ll i = 0;i <= m + n;i++)
		cout << (ll)(f[i].real() / len + 0.5) << " ";
	return 0;
}
\end{lstlisting}

\subsection{快速数论变换(NTT)}
\begin{lstlisting}[caption={}]
static constexpr int N = 4e6 + 6,g = 3,mod = 998244353;
i64 ni,gi;
i64 qpow(i64 a,i64 b){
    i64 ans = 1;
    for(;b;a = a * a % mod,b >>= 1)
        if(b & 1) ans = ans * a % mod;
    return ans;
}
void NTT(vector<i64>& A,int n,int op){
    vector<i64> R(n + 1,0);
	for(int i = 0;i < n;i++)
		R[i] = R[i / 2] / 2 + (i & 1) * (n / 2);
	for(int i = 0;i < n;i++)
		if(i < R[i]) swap(A[i],A[R[i]]);
	for(int i = 2;i <= n;i <<= 1){
        i64 g1 = qpow(op == 1 ? g : gi,(mod - 1) / i);
        for(int j = 0;j < n;j += i){
            i64 gk = 1;
            for(int k = j;k < j + i / 2;k++){
                i64 x = A[k],y = gk * A[k + i / 2] % mod;
                A[k] = (x + y) % mod;
                A[k + i / 2] = (x - y + mod) % mod;
                gk = gk * g1 % mod;
            }
        }
    }
}
void clac(vector<i64>& a,vector<i64>& b){
    const int la = a.size(),lb = b.size();
	int len = 1 << max((int)ceil(log2(la + lb)),1);
    vector<i64> f(3 * len + 1,0),h(3 * len + 1,0);
	for(int i = 0;i < la;i++) f[i] = a[i];
	for(int i = 0;i < lb;i++) h[i] = b[i];
    gi = qpow(g,mod - 2),ni = qpow(len,mod - 2);
	NTT(f,len,1),NTT(h,len,1);
	for(int i = 0;i <= len;i++) 
        f[i] = (f[i] * h[i]) % mod;
	NTT(f,len,-1);
	for(int i = 0;i <= la + lb - 2;i++) 
        cout << (f[i] * ni) % mod << " ";
}
\end{lstlisting}

\subsection{高斯消元}
\begin{lstlisting}[caption={}]
static constexpr double eps = 1e-8;
int Gauss(std::vector<std::vector<double>>& a){
	int n = a.size() - 1,now = 1;
	for(int i = 1;i <= n;i++){ //枚举列
		int r = now;
		for(int j = i;j <= n;j++){ //找最大减小误差
			if(fabs(a[j][i]) > fabs(a[r][i])) r = j;
		}
		if(fabs(a[r][i]) < eps) continue;
		if(r != now) swap(a[now],a[r]);
		for(int j = 1;j <= n;j++){
			if(j == now) continue;
			double div = a[j][i] / a[now][i];
			for(int k = i;k <= n + 1;k++) a[j][k] -= a[now][k] * div;
		}
		now++;
	}
	if(now <= n){
		while(now <= n){
			if(fabs(a[now][n + 1]) > eps) return -1;
			now++;
		}
		return 0;
	}
	for(int i = 1;i <= n;i++) a[i][n + 1] /= a[i][i];
	return 1;
} // 1 有解,-1无解,0无穷解
\end{lstlisting}

\subsection{rho算法与MR素性测试}
\begin{lstlisting}[caption={}]
std::mt19937_64 rng {std::chrono::steady_clock::now().time_since_epoch().count()};
i64 gcd(i64 a, i64 b);
i64 qpow(i64 a, i64 n, i64 p);
bool Mii64er_Rabin(i64 x){
    if (x < 3) return x == 2;
    if (x % 2 == 0) return false;
    constexpr i64 A[] = {2, 325, 9375, 28178, 450775, 9780504, 1795265022};
	i64 d = x - 1, r = 0;
    while (d % 2 == 0) d /= 2, ++r;
    for (auto a : A){
        i64 v = qpow(a, d, x);
        if(v <= 1 || v == x - 1) continue;
        for(int i = 0; i < r; ++i){
            v = (__int128)v * v % x;
            if(v == x - 1 && i != r - 1){
                v = 1;
                break;
            }
            if(v == 1)  return false;
        }
        if(v != 1) return false;
    }
    return true;
}
i64 Poi64ard_Rho(i64 N){
    if (N == 4) return 2;
    if (Mii64er_Rabin(N))return N;
    while(true){
        i64 c = rng() % (N - 1) + 1;
        auto f = [=](i64 x) {return ((__int128)x * x + c) % N;};
        i64 t = 0, r = 0, p = 1,q;
        do{
            for(int i = 0; i < 128; ++i){
                t = f(t),r = f(f(r));
                if (t == r || (q = (__int128)p * abs(t - r) % N) == 0)
                    break;
                p = q;
            }
            i64 d = gcd(p, N);
            if (d > 1) return d;
        }while (t != r);
    }
	return N;
}
unordered_map<i64, i64> um;
i64 max_prime_factor(i64 x){
    if(um.count(x)) return um[x];
    i64 fac = Pollard_Rho(x);
    if(fac == 1 || fac == x) um[x] = x;
    else um[x] = max(max_prime_factor(fac),max_prime_factor(x / fac));
    return um[x];
}
\end{lstlisting}

\subsection{求解连续按位异或}
以 $\mathcal O(1)$ 复杂度计算 $0\oplus1\oplus\dots\oplus n$ 。
\begin{lstlisting}[caption={}]
unsigned xor_n(unsigned n) {
    unsigned t = n & 3;
    if (t & 1) return t / 2u ^ 1;
    return t / 2u ^ n;
}
\end{lstlisting}
\begin{lstlisting}[caption={}]
i64 xor_n(i64 n) {
    if (n % 4 == 1) return 1;
    else if (n % 4 == 2) return n + 1;
    else if (n % 4 == 3) return 0;
    else return n;
}
\end{lstlisting}

\subsection{其他结论}

\begin{itemize}
    \item 欧拉降幂公式(扩展欧拉定理):
		$$a^b\equiv
        \begin{cases}
        a^{b\bmod\varphi(m)},\,&\gcd(a,\,m)=1\\
        a^b,&\gcd(a,\,m)\ne1,\,b<\varphi(m)\\
        a^{b\bmod\varphi(m)+\varphi(m)},&\gcd(a,\,m)\ne1,\,b\ge\varphi(m)
        \end{cases}$$
    \item 约数个数定理: 设$n = p_1^{a_1}\times p_1^{a_1} \times \cdots \times p_m^{a_m}$,其中$p_1,p_2,\cdots,p_3$均为互不相等的质数，则$n$的因数个数为：
$d(n) = \prod\limits_{i = 1}^{m}(a_i + 1)$.

    \item 约数和定理: $\sigma(n) = \sum\limits_{d | n} d = \prod\limits_{i = 1}^{k} \frac{p_i^{a_i+1} - 1}{p_i - 1}$
    \item 约数个数前缀和: $\sum \limits _{i = 1}^{n} d(i) = \sum \limits _{i = 1}^{n}{\lfloor \frac{n}{i} \rfloor}$
    \item 约数和前缀和: $\sum \limits _{i = 1}^{n} \sigma(i) = \sum \limits _{i = 1}^{n}{\lfloor \frac{n}{i} \rfloor i}$
    \item 勒让德定理（Legendre 定理）: 
        对于质数$p$，函数$v_p(n)$为n的标准分解后$p$的次数，有
        $$v_p(n!) = \sum\limits_{i = 1}^{\infty} \lfloor \frac{n}{p^i} \rfloor$$

        令函数$s_p(n)$为$n$在$p$进制下的数位和有：
        $$v_p(n!) = \frac{n - s_p(n)}{p - 1}$$
        $p = 2$时有：
        $$v_2(n!) = n - popcount(n)$$

        kummer定理,二项式系数:
        $$v_p((^n_m)) = \frac{s_p(m)+s_p(n - m)-s_p(n)}{p - 1}$$
    
    \item 霍尔定理(Hall 定理) :
            二分图存在完美匹配的判定：对于一个二分图$G(V_1,V_2,E) (|V_1| \le |V_2|)$，对与$\forall v_1 \subseteq V_1,v2 = \{v_j | (v_i,v_j) \in E,v_i \in V_1,v_j \in V_2 \}$，其存在完美匹配的条件为$\forall v_1 
            \in V_1,|v_1| <= |v_2|$.
    \item Dilworth定理: 偏序集能划分成的最少的全序集个数等于最大反链的元素个数，即数组的不重叠的上升子序列数量最小值 = 逆序数组的最长上升子序列长度。
    \item 柯尼希定理(Kőnig 定理): 二分图最小点覆盖的点数=最大匹配数。
    \item 切比雪夫距离和曼哈顿距离:
        将一个点$(x,y)$的坐标变为$(x+y,x - y)$后,原坐标系中的曼哈顿距离 = 新坐标系中的切比雪夫距离

        将一个点$(x,y)$的坐标变为$(\frac{x + y}{2},\frac{x - y}{2})$ 后,原坐标系中的切比雪夫距离 = 新坐标系中的曼哈顿距离
    \item 基姆拉尔森公式:用于计算某一天是星期几:
        $$h = (d + \lfloor \frac{13(m + 1)}{5} \rfloor + y + \lfloor \frac{y}{4} \rfloor + \lfloor \frac{c}{4} \rfloor + 5c) \mod 7$$ 
    \item $ a - b <= a \oplus b <= a + b  $
    \item 对于给定的 $X$ 和序列 $[a_1,a_2,…,a_n]$ ，有：${X=(X \&a_1)or(X\&a_2)or…or(X\&a_n)}$。
    \item $gcd(x,y) = x \oplus y \to gcd(x,y)=y-x$
    \item $gcd(x,y) = x + y \to gcd(x,y)=max(x,y)$
    \item 以下式子成立： $\gcd (a, m) = \gcd(a+x,m) \Leftrightarrow  \gcd(a, m)=\gcd(x,m)$ 。
    求解上式满足条件的 $x$ 的数量即为求比 $\dfrac{m}{\gcd(a,m)}$ 小且与其互质的数的个数，即用欧拉函数求解 $\varphi \Big(\dfrac{m}{\gcd(a,m)} \Big)$ 。
    \item 令 $\left\{\begin{matrix} X=a+b\\ 
            Y=a\oplus b
            \end{matrix}\right.$ ，如果该式子有解，那么存在前提条件 $\left\{\begin{matrix} X \ge Y \\ 
            X,Y \text{同奇偶}
            \end{matrix}\right.$ ；进一步，此时最小的 $a$ 的取值为 $\dfrac{X-Y}{2}$
    \item $x+y=x|y+x\&y$ ，对于两个数字 $x$ 和 $y$ ，如果将 $x$ 变为 $x|y$ ，同时将 $y$ 变为 $x\&y$ ，
            那么在本质上即将 $x$ 二进制模式下的全部 $1$ 移动到了 $y$ 的对应的位置上
    \item 一个正整数 $x$ 异或、加上另一个正整数 $y$ 后奇偶性不发生变化：$a+b\equiv a\oplus b(\bmod2)$
    \item 贝叶斯公式：$$P(A_i|B) = \frac{P(A_iB)}{P(B)}=\frac{P(A_i)P(B|A_i)}{\displaystyle\sum^{n}_{j=1}P(A_j)P(B|A_j)}$$
    \item 全概率公式: $$P(B) = \sum\limits_{i=1}^{n}P(A_i)P(B | A_i)$$
    \item 全期望公式: $$E(Y) = E(E(Y|X))=\sum \limits_{X}E(Y|X)P(X)$$
    \item $\forall a,b,c \in  \mathbb{Z},\left \lfloor \frac{a}{bc}  \right \rfloor 
        = \left \lfloor \frac{\left \lfloor \frac{a}{b}  \right \rfloor }{c}  \right \rfloor $
\end{itemize}